% paper/outline.tex
% Detailed outline of the paper structure. Companion to
% paper/guiding_principles.md. Intended as a planning document, not for
% inclusion in the final paper. Not compiled standalone.

\documentclass{article}
\usepackage{amsmath,amssymb}
\usepackage{enumitem}
\usepackage[margin=1in]{geometry}

\title{Paper Outline: Local Income Taxes and Residential Sorting}
\author{}
\date{}

\begin{document}
\maketitle

\section{Introduction}
\begin{itemize}
  \item The question: do local income tax rate changes induce residential
    sorting?
  \item The institutional opportunity: PA EIT's piecewise rule + Act 32
    (2012) consolidation
  \item Headline findings preview (per-pp dose-response on Inc-bind cell,
    three estimating units)
  \item Contribution to the continuous stacked DiD literature
    (Cengiz 2019; Callaway, Goodman-Bacon \& Sant'Anna 2024;
    de Chaisemartin \& D'Haultf{\oe}uille 2024)
\end{itemize}

\section{Institutional Background}
\begin{itemize}
  \item PA EIT structure: muni resident rate + SD resident rate (home
    stack) vs.\ work-muni nonresident rate
  \item Piecewise rule:
    \( T = \max(\tau^M + \tau^{SD},\ \tau^{nr}) \)
  \item DCED rate-setting and reporting
  \item Act 32 of 2008 (effective 2012): consolidated employer
    withholding; the natural experiment for binding
\end{itemize}

\section{Conceptual Framework}
\begin{itemize}
  \item Why piecewise creates a within-design placebo (Inc-nobind)
  \item Cell decomposition:
    Inc-bind / Inc-becomes / Inc-nobind /
    Dec-bind / Dec-becomes-non / Dec-nobind
  \item Predicted-null cell: Inc-nobind
  \item Predicted-magnitude ordering across cells
\end{itemize}

\section{Data}

\subsection{LEHD LODES Origin--Destination at the block level}
\begin{itemize}
  \item Coverage 2002--2019; PA only
  \item Block-level OD; aggregated up where required
\end{itemize}

\subsection{Tax rate panel}
\begin{itemize}
  \item DCED muni and SD rate panels
  \item Construction of total\_eit, home\_stack, muni\_nonres\_eit
  \item Coverage and event counts
\end{itemize}

\subsection{Municipal fiscal controls}
\begin{itemize}
  \item AFR (Annual Financial Report) muni share-based controls:
    share\_police, share\_streets, share\_debt\_svc,
    share\_eit, share\_re, share\_intergov
\end{itemize}

\subsection{School district fiscal controls}
\begin{itemize}
  \item Census Bureau F-33 SD revenue/expenditure data; share-based
    controls parallel to muni AFR shares: share\_state\_rev\_sd,
    share\_local\_rev\_sd, share\_property\_tax\_sd,
    share\_instructional\_sd, share\_support\_sd, share\_capital\_sd,
    share\_debt\_int\_sd
\end{itemize}

\subsection{Three estimating panels}

\paragraph{Municipality-level panel}
Unit = cousub $\times$ year. Workplace and residence outcomes
constructed by aggregating LODES OD over blocks within the cousub.
Treatment side: muni resident rate is exact; school district resident
rate is assigned as max(sd\_eit) across the school districts the muni
reports to in DCED. Justification: 97.6\% of munis report to one SD;
the max is correct for those, conservative for the 2.4\% multi-SD
munis. Fiscal controls: muni AFR shares only.

\paragraph{School-district-level panel}
Unit = SD $\times$ year. Outcomes constructed by aggregating LODES OD
over blocks within the SD. Treatment side: SD resident rate is exact;
municipal resident rate is assigned as the within-SD aggregation across
the SD's constituent munis. Fiscal controls: F-33 SD shares only.

\paragraph{Block-level panel}
Unit = block $\times$ year. Each block is assigned its own muni and its
own SD. Both muni and SD resident rates enter the regression at their
native level with no within-jurisdiction aggregation. Outcome variables
are block-defined: log\_home\_count, log\_work\_count,
home\_in\_muni\_share, share\_work\_outside. Fiscal controls: both muni
AFR shares and F-33 SD shares (each block is exposed to both
governments).

\subsection{Sample restrictions}
\begin{itemize}
  \item 2002--2019 (drops COVID-affected and ARPA-affected fiscal years)
  \item Min thresholds (e.g., total\_jobs $> 50$ at cousub level) where
    applicable
\end{itemize}

\section{Empirical Strategy}

\subsection{Continuous stacked DiD: main framework}
\begin{itemize}
  \item Stacked cohort design (Cengiz)
  \item Treatment intensity:
    \(\Delta\tau_{it} = \tau_{it} - \tau_{i,c-1}\), time-varying, zero
    pre-event
  \item NC controls (never-changed jurisdictions)
\end{itemize}

\subsection{Cell-decomposed dose-response specification}
\begin{equation*}
\begin{split}
Y_{it} = \alpha_{i \times s} + \gamma_{t \times s}
&+ \sum_{c \neq NC} \beta^{M}_c \mathbf{1}\{\text{cell}^M_i = c\}
    \cdot \Delta\tau^M_{it} \\
&+ \sum_{c \neq NC} \beta^{SD}_c \mathbf{1}\{\text{cell}^{SD}_i = c\}
    \cdot \Delta\tau^{SD}_{it} \\
&+ X_{it}\phi + \varepsilon_{it}
\end{split}
\end{equation*}
Run at all three estimating units: muni, SD, block.

\subsection{Fixed effects}
Unit-by-stack and year-by-stack. Cengiz-style sub-experiment
isolation.

\subsection{Clustering}
Cluster at the policy-jurisdiction level. Muni: $\sim$muni. SD:
$\sim$sd\_id. Block: two-way $\sim$muni + sd\_id.

\subsection{Pre-trend diagnostic: realized-dose event study}
\begin{itemize}
  \item Spec: $i(\text{rel}, \widetilde{D}_m, \text{ref} = -1)$, with
    $\widetilde{D}_m$ defined as the unit's eventual cumulative
    $\Delta\tau$ at $+5$
  \item Pre-event leads identified from cross-unit variation in
    eventual treatment intensity
  \item Reported for each estimating unit and each headline outcome
\end{itemize}

\section{Main Results}

\subsection{Municipality-level results}
All cells $\times$ all outcomes. Highlight: Inc-bind, Inc-becomes,
Inc-nobind (placebo).

\subsection{School-district-level results}
Same structure. Discuss SD-side framing.

\subsection{Block-level results}
Block-defined outcomes. Note: treatment variation now exact (no
aggregation function).

\subsection{Cross-level comparison}
Triangulation table: same cell $\times$ outcome at all three levels.

\section{Heterogeneity Decompositions}

\subsection{Pre-/post-Act 32 split}
Run for each estimating unit (muni, SD, block). Predicted: Inc-bind
effects concentrated post-2012 (when EIT became binding via
consolidated withholding). Headline cell $\times$ pre-Act 32 vs
post-Act 32 interactions.

\subsection{Industry decomposition (job-side outcomes)}
Services vs non-services using LODES WAC supersector codes. Run on
log\_total\_jobs (or log\_work\_count at block level). Predicted:
services more elastic than non-services. Run for each estimating unit.

\subsection{Muni-type heterogeneity}
City vs borough vs township. Run on each level. Predicted: cities
(denser, more exit options) show larger effects.

\section{Robustness}

\subsection{Realized-dose event-study (pre-trend test)}
Per cell $\times$ per outcome $\times$ per estimating unit. Pass/fail
summary table. ES figures for Inc-bind on each headline outcome.

\subsection{Alternative control definitions}
Tighter ``no events in 8-year window'' controls (Cengiz-style).
Compare to NC-only headline.

\subsection{Sample restrictions}
Drop flip-flopper munis. Drop top-quartile pre-event growth munis
(gentrification check).

\subsection{Single-event subset}
Munis with one event only.

\section{Discussion}
\begin{itemize}
  \item Magnitude in welfare / fiscal-calibration terms
  \item Implications for local government fiscal autonomy
  \item Limits: PA-only, residence margin only, etc.
\end{itemize}

\section{Conclusion}

\section*{What's dropped from the previous paper structure}
\begin{itemize}
  \item Single-event TWFE as a separate estimator section
  \item DCDH cell-by-cell as a separate estimator
  \item Stacked binary DiD as a separate estimator
  \item Piecewise binding model on cousub-pair flows
  \item Resident vs nonresident-side parallel build (nonresident deferred)
\end{itemize}

\section*{What's added}
\begin{itemize}
  \item Three-level estimation structure (muni / SD / block) with
    explicit treatment-aggregation rules
  \item Industry decomposition (services vs non-services) on job-side
    outcomes
  \item Pre/post Act 32 split as a heterogeneity section, run at each level
  \item Realized-dose event study as the canonical pre-trend diagnostic
\end{itemize}

\end{document}
